% BHU PYQ -- Manually transcribed & error-corrected
% Paper: BSC-04A: Elements of Geology - II (Ancillary)
% Exam: B.Sc. (Hons.) Semester IV Examination, 2013-14

\documentclass[11pt,a4paper]{article}
\usepackage[a4paper,top=2.5cm,bottom=2.5cm,left=2.8cm,right=2.8cm,headheight=14pt]{geometry}
\usepackage{amsmath,amssymb}
\usepackage[T1]{fontenc}
\usepackage[utf8]{inputenc}
\usepackage{lmodern,microtype,fancyhdr,enumitem,hyperref,lastpage,parskip}

\hypersetup{colorlinks=true,urlcolor=black,linkcolor=black,
  pdftitle={BSC-04A Elements of Geology - II (Ancillary) -- BHU B.Sc. Sem IV 2013-14}}
\pagestyle{fancy}\fancyhf{}
\renewcommand{\headrulewidth}{0.4pt}\renewcommand{\footrulewidth}{0.4pt}
\fancyhead[L]{\small   \textit{Banaras Hindu University}}
\fancyhead[R]{\small   \textit{BSC-04A: Elements of Geology - II (Ancillary)}}
\fancyfoot[C]{\small Page  \thepage\ of \pageref{LastPage}}


\newlist{parts}{enumerate}{2}
\setlist[parts,1]{label=(\alph*),leftmargin=*,itemsep=5pt,topsep=3pt}

\newlist{subparts}{enumerate}{2}
\setlist[subparts,1]{label=(\roman*),leftmargin=*,itemsep=3pt,topsep=2pt}

\newcommand{\pts}[1]{\hfill   \textit{[#1]}}


\begin{document}

\begin{center}
  {\large\bfseries BANARAS HINDU UNIVERSITY}\\[2pt]
  {\normalsize B.Sc. (Hons.) --- Semester IV Examination, 2013-14}\\[6pt]
  \rule{0.8\linewidth}{0.5pt}\\[6pt]
  {\large\bfseries Paper: BSC-04A --- Elements of Geology - II (Ancillary)}\\[4pt]
  {\normalsize Time: 3 Hours \hfill Maximum Marks: 70}\\[2pt]
  {\small Ref. No.: BSC04A Ancillary}
\end{center}
\rule{\linewidth}{0.4pt}
\smallskip



\noindent   \textit{Instructions: Answer   \textbf{five} questions in all, selecting at least   \textbf{two} each from Sections A and B, and Section-C which is compulsory. The figures in the right-hand margin indicate marks.}
\bigskip

\bigskip


\noindent {\large\bfseries SECTION --- A}
\rule{\linewidth}{0.4pt}\smallskip




\noindent   \textbf{Question 1.} Explain how Palaeontology supports the phenomenon of ``Organic Evolution''. \pts{14}

\medskip


\noindent   \textbf{Question 2.} Discuss petrifaction and its kinds. What change(s) occur(s) with respect to the original organism? \pts{14}

\medskip


\noindent   \textbf{Question 3.} Write explanatory notes on   \textbf{any four} of the following: \pts{$4  \times 3\frac{1}{2} = 14$}

\begin{parts}
  \item Subdisciplines of Palaeontology
  \item Distinction between a true shell and a cast
  \item The process of carbonization. What deformation does the fossil undergo with respect to the actual organism?
  \item What are the TWO major requirements for fossilization?
  \item Basic differences between `Coprolite' and `Gastrolith'.
\end{parts}

\bigskip


\noindent {\large\bfseries SECTION --- B}
\rule{\linewidth}{0.4pt}\smallskip




\noindent   \textbf{Question 4.} Discuss the physical divisions of India with a neat diagram. \pts{14}

\medskip


\noindent   \textbf{Question 5.} What is the Geological Time Scale? Write a note on how it was developed and its various uses. \pts{14}

\medskip


\noindent   \textbf{Question 6.} Define mass extinction. How many mass extinctions are generally believed to have occurred on the Earth and during which geological time periods? Add a note on the possible causes of mass extinction. \pts{14}

\bigskip


\noindent {\large\bfseries SECTION --- C}
\rule{\linewidth}{0.4pt}\smallskip




\noindent   \textbf{Question 7.} Answer the following questions: \pts{14}

\begin{parts}
  \item Explain the etymology of Palaeontology.
  \item What is the oldest trace fossil?
  \item Which of the following is/are `true' Body Fossil(s)?
  
\begin{subparts}
    \item Mesozoic dinosaur footprints
    \item Chemically preserved `Mummies' of Pharaohs
    \item Oligocene insects preserved in amber
    \item Pleistocene fossil `Mammoths' of Tundra
  \end{subparts}
  \item Define a ``Steinkern''.
  \item Name a `metallic' mineral deposit discovered through the study of fossils.
  \item Tarai-Bhabhar
  \item The principle of superposition
  \item Uniformitarianism
  \item Name the important Upper Gondwana plant fossils.
  \item Barren Measures
\end{parts}

\vfill\rule{\linewidth}{0.4pt}

\begin{center}\small
\href{https://bhu-exam.vercel.app/}{Click here to visit our website}\\[4pt]
\href{https://me-aryan.vercel.app/}{Click here to visit our contributor}\\[4pt]
\href{https://quashan-me.vercel.app/}{Click here to visit our contributor}
\end{center}
\end{document}
